\chapter{Kommutatoren und Unschärferelation}

\section{Ziel dieses Kapitels}

In diesem Kapitel geht es um eine der wichtigsten Strukturen der Quantenmechanik: Operatoren können im Allgemeinen nicht vertauscht werden. In der klassischen Mechanik ist für zwei gewöhnliche Zahlen oder Funktionen meist klar, dass
\[
    ab=ba
\]
gilt. In der Quantenmechanik wirken Observablen aber als Operatoren auf Zustände. Dann kann die Reihenfolge der Anwendung physikalisch relevant sein:
\[
    AB\psi \neq BA\psi.
\]
Der Unterschied zwischen beiden Reihenfolgen wird durch den Kommutator gemessen.

\begin{conceptbox}
Kommutatoren sind nicht nur ein Rechentrick. Sie sagen aus, ob zwei Observablen gleichzeitig scharf messbar sind, ob eine Größe eine Erhaltungsgröße ist und welche Unschärferelation zwischen zwei Observablen gilt.
\end{conceptbox}

Die zentralen Fragen dieses Kapitels sind:
\begin{itemize}
    \item Was ist ein Kommutator und wie rechnet man mit ihm?
    \item Warum gilt für Ort und Impuls die fundamentale Relation $[X_j,P_k]=i\hbar\delta_{jk}$?
    \item Wie berechnet man Kommutatoren mit Funktionen von Operatoren?
    \item Wie hängen Kommutatoren mit simultaner Messbarkeit zusammen?
    \item Wie folgt die Heisenbergsche Unschärferelation aus der Hilbertraumstruktur?
    \item Welche Zustände minimieren die Orts-Impuls-Unschärfe?
    \item Wie entstehen Ehrenfest-Gleichungen aus Kommutatoren mit dem Hamiltonoperator?
\end{itemize}

\section{Definition des Kommutators}

\begin{definitionbox}
Für zwei lineare Operatoren $A$ und $B$ ist der Kommutator definiert durch
\[
    [A,B] := AB-BA.
\]
Man sagt, $A$ und $B$ kommutieren, wenn
\[
    [A,B]=0
\]
gilt.
\end{definitionbox}

Der Kommutator misst also, ob die Reihenfolge der Operatoren egal ist. Wirkt man auf einen Zustand $\psi$, so ist
\[
    [A,B]\psi = A(B\psi)-B(A\psi).
\]
Wenn dieser Ausdruck für alle Zustände verschwindet, dann sind $A$ und $B$ vertauschbar.

\begin{notebox}
Die Gleichung $[A,B]=0$ bedeutet nicht nur, dass bei einer einzelnen Rechnung zufällig dasselbe herauskommt. Sie bedeutet, dass die beiden Operatoren als Operatoren identisch vertauschbar sind, also auf ihrem gemeinsamen Definitionsbereich dieselbe Wirkung haben.
\end{notebox}

\subsection{Elementare Eigenschaften}

Kommutatoren erfüllen einige Rechenregeln, die man ständig verwendet.

\begin{formulabox}
Für Operatoren $A,B,C$ und Zahlen $\alpha,\beta\in\mathbb C$ gilt:
\[
    [A,B]=-[B,A],
\]
\[
    [A,\alpha B+\beta C]=\alpha[A,B]+\beta[A,C],
\]
\[
    [\alpha A+\beta B,C]=\alpha[A,C]+\beta[B,C].
\]
Außerdem gelten die Produktregeln
\[
    [A,BC]=[A,B]C+B[A,C],
\]
\[
    [AB,C]=A[B,C]+[A,C]B.
\]
\end{formulabox}

\begin{derivationbox}
Die Produktregel folgt direkt aus der Definition:
\[
\begin{aligned}
    [A,BC]
    &=ABC-BCA \\
    &=ABC-BAC+BAC-BCA \\
    &=(AB-BA)C+B(AC-CA) \\
    &=[A,B]C+B[A,C].
\end{aligned}
\]
Analog erhält man
\[
    [AB,C]=A[B,C]+[A,C]B.
\]
\end{derivationbox}

\subsection{Kommutator und Antikommutator}

Neben dem Kommutator verwendet man auch den Antikommutator.

\begin{definitionbox}
Der Antikommutator zweier Operatoren $A$ und $B$ ist
\[
    \{A,B\}:=AB+BA.
\]
\end{definitionbox}

Jedes Produkt zweier Operatoren kann in einen symmetrischen und einen antisymmetrischen Anteil zerlegt werden:
\[
    AB=\frac12\{A,B\}+\frac12[A,B].
\]
Der Antikommutator ist besonders wichtig in der genauen Form der Robertson-Schrödinger-Unschärferelation. Für die einfache Heisenberg-Form reicht oft der Kommutatoranteil.

\section{Der fundamentale Heisenberg-Kommutator}

\subsection{Orts- und Impulsoperator in Ortsdarstellung}

In Ortsdarstellung wirken Orts- und Impulsoperator auf eine Wellenfunktion $\psi(\vec x)$ durch
\[
    (X_j\psi)(\vec x)=x_j\psi(\vec x),
    \qquad
    (P_j\psi)(\vec x)=-i\hbar\frac{\partial}{\partial x_j}\psi(\vec x).
\]
Für den Kommutator $[X_j,P_k]$ müssen wir die beiden Reihenfolgen vergleichen.

\begin{derivationbox}
Wir lassen $[X_j,P_k]$ auf eine Testfunktion $\psi(\vec x)$ wirken:
\[
\begin{aligned}
    (X_jP_k\psi)(\vec x)
    &=x_j\left(-i\hbar\frac{\partial \psi}{\partial x_k}\right), \\
    (P_kX_j\psi)(\vec x)
    &=-i\hbar\frac{\partial}{\partial x_k}\left(x_j\psi(\vec x)\right) \\
    &=-i\hbar\left(\delta_{jk}\psi(\vec x)+x_j\frac{\partial\psi}{\partial x_k}\right).
\end{aligned}
\]
Daher ist
\[
\begin{aligned}
    ([X_j,P_k]\psi)(\vec x)
    &=x_j\left(-i\hbar\frac{\partial\psi}{\partial x_k}\right)
    +i\hbar\left(\delta_{jk}\psi+x_j\frac{\partial\psi}{\partial x_k}\right) \\
    &=i\hbar\delta_{jk}\psi(\vec x).
\end{aligned}
\]
Da dies für alle geeigneten $\psi$ gilt, folgt
\[
    [X_j,P_k]=i\hbar\delta_{jk}\,\mathbb{1}.
\]
Wenn die Identität klar ist, schreibt man kurz
\[
    [X_j,P_k]=i\hbar\delta_{jk}.
\]
\end{derivationbox}

Da verschiedene Ortskomponenten nur Multiplikation mit verschiedenen Koordinaten sind, gilt
\[
    [X_j,X_k]=0.
\]
Da partielle Ableitungen nach verschiedenen Koordinaten miteinander vertauschen, gilt außerdem
\[
    [P_j,P_k]=0.
\]

\begin{formulabox}
Die kanonischen Kommutatorrelationen lauten
\[
    [X_j,X_k]=0,
    \qquad
    [P_j,P_k]=0,
    \qquad
    [X_j,P_k]=i\hbar\delta_{jk}.
\]
Sie bilden die mathematische Grundstruktur der nichtrelativistischen Quantenmechanik.
\end{formulabox}

\begin{notebox}
Manchmal steht in Aufgaben $P_j=\frac{\hbar}{i}\partial_{x_j}$. Das ist dasselbe wie $P_j=-i\hbar\partial_{x_j}$, weil $1/i=-i$ gilt.
\end{notebox}

\section{Kommutatoren mit Potenzen und Funktionen}

\subsection{Kommutator mit einer Potenz des Impulses}

Aus
\[
    [X_j,P_j]=i\hbar
\]
folgt eine ganze Familie von Kommutatorrelationen. Besonders wichtig ist
\[
    [X_j,P_j^n]=i\hbar nP_j^{n-1}.
\]

\begin{derivationbox}
Wir verwenden die Produktregel
\[
    [A,BC]=[A,B]C+B[A,C].
\]
Für $n=2$ gilt
\[
\begin{aligned}
    [X_j,P_j^2]
    &=[X_j,P_jP_j] \\
    &=[X_j,P_j]P_j+P_j[X_j,P_j] \\
    &=i\hbar P_j+P_j i\hbar \\
    &=2i\hbar P_j.
\end{aligned}
\]
Für allgemeines $n$ erhält man durch Induktion
\[
\begin{aligned}
    [X_j,P_j^n]
    &=[X_j,P_jP_j^{n-1}] \\
    &=[X_j,P_j]P_j^{n-1}+P_j[X_j,P_j^{n-1}] \\
    &=i\hbar P_j^{n-1}+P_j\,i\hbar(n-1)P_j^{n-2} \\
    &=i\hbar nP_j^{n-1}.
\end{aligned}
\]
\end{derivationbox}

Analog gilt
\[
    [P_j,X_j^n]=-i\hbar nX_j^{n-1}.
\]

\subsection{Kommutatoren mit Funktionen}

Wenn eine Funktion als Potenzreihe geschrieben werden kann,
\[
    G(P_j)=\sum_{n=0}^{\infty}a_nP_j^n,
\]
dann folgt aus Linearität des Kommutators
\[
\begin{aligned}
    [X_j,G(P_j)]
    &=\sum_{n=0}^{\infty}a_n[X_j,P_j^n] \\
    &=\sum_{n=1}^{\infty}a_n i\hbar nP_j^{n-1} \\
    &=i\hbar\frac{\partial G}{\partial P_j}.
\end{aligned}
\]
Genauso gilt für eine Funktion $F(X_j)$:
\[
    [P_j,F(X_j)]=-i\hbar\frac{\partial F}{\partial X_j}.
\]

\begin{formulabox}
Für hinreichend glatte beziehungsweise als Potenzreihe darstellbare Funktionen gilt
\[
    [X_j,G(P_j)]=i\hbar\frac{\partial G}{\partial P_j},
    \qquad
    [P_j,F(X_j)]=-i\hbar\frac{\partial F}{\partial X_j}.
\]
Insbesondere gilt in mehreren Dimensionen
\[
    [P_j,V(\vec X)]=-i\hbar\frac{\partial V}{\partial x_j}(\vec X).
\]
\end{formulabox}

\begin{examplebox}
Für $G(P)=P^2/(2m)$ erhält man
\[
    [X,G(P)]=i\hbar\frac{P}{m}.
\]
Für $F(X)=V(X)$ erhält man
\[
    [P,V(X)]=-i\hbar V'(X).
\]
Diese beiden Relationen sind genau die Bausteine des Ehrenfest-Theorems.
\end{examplebox}

\section{Hamiltonoperator, Erhaltungsgrößen und Ehrenfest-Theorem}

\subsection{Zeitentwicklung eines Erwartungswertes}

Sei $A$ eine Observable. Ihr Erwartungswert im Zustand $\psi(t)$ ist
\[
    \langle A\rangle_t=\langle\psi(t),A\psi(t)\rangle.
\]
Wenn $A$ eventuell explizit zeitabhängig ist, dann gilt
\[
    \frac{\dd}{\dd t}\langle A\rangle_t
    =\left\langle\frac{\partial A}{\partial t}\right\rangle_t
    +\frac{i}{\hbar}\langle[H,A]\rangle_t.
\]

\begin{derivationbox}
Aus der Schrödingergleichung
\[
    i\hbar\frac{\partial}{\partial t}\psi=H\psi
\]
folgt
\[
    \dot\psi=-\frac{i}{\hbar}H\psi.
\]
Für den Bra gilt entsprechend
\[
    \frac{\dd}{\dd t}\langle\psi|=\frac{i}{\hbar}\langle\psi|H,
\]
sofern $H=H^\dagger$ gilt. Damit ist
\[
\begin{aligned}
    \frac{\dd}{\dd t}\langle A\rangle
    &=\frac{\dd}{\dd t}\langle\psi|A|\psi\rangle \\
    &=\langle\dot\psi|A|\psi\rangle
      +\left\langle\psi\left|\frac{\partial A}{\partial t}\right|\psi\right\rangle
      +\langle\psi|A|\dot\psi\rangle \\
    &=\frac{i}{\hbar}\langle\psi|HA|\psi\rangle
      +\left\langle\frac{\partial A}{\partial t}\right\rangle
      -\frac{i}{\hbar}\langle\psi|AH|\psi\rangle \\
    &=\left\langle\frac{\partial A}{\partial t}\right\rangle
      +\frac{i}{\hbar}\langle[H,A]\rangle.
\end{aligned}
\]
\end{derivationbox}

\begin{formulabox}
Eine Observable $A$ ohne explizite Zeitabhängigkeit ist eine Erhaltungsgröße, wenn
\[
    [H,A]=0
\]
gilt. Dann ist
\[
    \frac{\dd}{\dd t}\langle A\rangle_t=0.
\]
\end{formulabox}

\subsection{Ehrenfest-Theorem}

Für den Standard-Hamiltonoperator
\[
    H=\frac{\vec P^{\,2}}{2m}+V(\vec X)
\]
ergeben sich aus den Kommutatoren direkt die quantenmechanischen Bewegungsgleichungen für Erwartungswerte.

\begin{derivationbox}
Zuerst berechnen wir $[H,X_j]$:
\[
\begin{aligned}
    [H,X_j]
    &=\left[\sum_{k=1}^3\frac{P_k^2}{2m}+V(\vec X),X_j\right] \\
    &=\sum_{k=1}^3\frac{1}{2m}[P_k^2,X_j]+[V(\vec X),X_j].
\end{aligned}
\]
Da $V(\vec X)$ nur eine Funktion der Ortsoperatoren ist, gilt $[V(\vec X),X_j]=0$. Weiter ist
\[
    [P_k^2,X_j]=P_k[P_k,X_j]+[P_k,X_j]P_k.
\]
Aus $[X_j,P_k]=i\hbar\delta_{jk}$ folgt $[P_k,X_j]=-i\hbar\delta_{jk}$. Also
\[
    [P_k^2,X_j]=-2i\hbar\delta_{jk}P_k.
\]
Damit folgt
\[
    [H,X_j]=-\frac{i\hbar}{m}P_j.
\]
Somit ist
\[
    \frac{\dd}{\dd t}\langle X_j\rangle
    =\frac{i}{\hbar}\langle[H,X_j]\rangle
    =\frac{1}{m}\langle P_j\rangle.
\]
\end{derivationbox}

\begin{derivationbox}
Nun berechnen wir $[H,P_j]$:
\[
\begin{aligned}
    [H,P_j]
    &=\left[\sum_{k=1}^3\frac{P_k^2}{2m}+V(\vec X),P_j\right] \\
    &=\sum_{k=1}^3\frac{1}{2m}[P_k^2,P_j]+[V(\vec X),P_j].
\end{aligned}
\]
Da die Impulskomponenten miteinander kommutieren, ist der erste Term null. Außerdem gilt
\[
    [P_j,V(\vec X)]=-i\hbar\frac{\partial V}{\partial x_j},
\]
also
\[
    [V(\vec X),P_j]=i\hbar\frac{\partial V}{\partial x_j}.
\]
Damit folgt
\[
    \frac{\dd}{\dd t}\langle P_j\rangle
    =\frac{i}{\hbar}\langle[H,P_j]\rangle
    =-\left\langle\frac{\partial V}{\partial x_j}\right\rangle.
\]
\end{derivationbox}

\begin{formulabox}
Das Ehrenfest-Theorem lautet
\[
    \frac{\dd}{\dd t}\langle X_j\rangle=\frac{1}{m}\langle P_j\rangle,
    \qquad
    \frac{\dd}{\dd t}\langle P_j\rangle=-\left\langle\frac{\partial V}{\partial x_j}\right\rangle.
\]
Es ist die quantenmechanische Entsprechung der klassischen Gleichungen
\[
    \dot x_j=\frac{p_j}{m},
    \qquad
    \dot p_j=-\frac{\partial V}{\partial x_j}.
\]
\end{formulabox}

\begin{notebox}
Wichtig ist der Unterschied zwischen
\[
    -\left\langle\frac{\partial V}{\partial x_j}\right\rangle
    \quad\text{und}\quad
    -\frac{\partial V}{\partial x_j}(\langle X\rangle).
\]
Diese Ausdrücke sind im Allgemeinen nicht gleich. Sie werden näherungsweise gleich, wenn das Wellenpaket schmal ist oder das Potential höchstens quadratisch ist.
\end{notebox}

\section{Kommutatoren und simultane Messbarkeit}

\subsection{Gemeinsame Eigenzustände}

Zwei Observablen $A$ und $B$ sind gleichzeitig scharf messbar, wenn der Zustand gleichzeitig Eigenzustand beider Operatoren ist:
\[
    A\psi=a\psi,
    \qquad
    B\psi=b\psi.
\]
Dann sind beide Messwerte in diesem Zustand sicher bestimmt.

Wenn $A$ und $B$ kommutieren, können sie unter geeigneten Bedingungen gemeinsam diagonalisiert werden. Das bedeutet, es gibt eine Basis aus gemeinsamen Eigenzuständen.

\begin{conceptbox}
Kommutierende Observablen sind die mathematische Formulierung simultaner Schärfe. Nichtkommutierende Observablen besitzen im Allgemeinen keine gemeinsame vollständige Eigenbasis und können daher nicht in allen Zuständen gleichzeitig scharf sein.
\end{conceptbox}

\subsection{Warum Kommutatoren gemeinsame Eigenzustände erzwingen können}

Angenommen, $A$ und $B$ haben einen gemeinsamen Eigenzustand $\psi$:
\[
    A\psi=a\psi,
    \qquad
    B\psi=b\psi.
\]
Dann gilt auf diesem Zustand
\[
\begin{aligned}
    [A,B]\psi
    &=AB\psi-BA\psi \\
    &=A(b\psi)-B(a\psi) \\
    &=bA\psi-aB\psi \\
    &=ba\psi-ab\psi=0.
\end{aligned}
\]
Ein gemeinsamer Eigenzustand liegt also im Kern des Kommutators.

\begin{notebox}
Aus $[A,B]\psi=0$ für einen einzelnen Zustand folgt noch nicht automatisch $[A,B]=0$ als Operatorgleichung. Die Operatorgleichung ist stärker: Sie muss auf allen Zuständen gelten.
\end{notebox}

\subsection{Entartung und vollständiger Satz kommutierender Observablen}

Wenn $A$ entartete Eigenwerte besitzt, kann es innerhalb eines entarteten Eigenraums mehrere unabhängige Eigenzustände geben. Dann reicht $A$ allein nicht aus, um den Zustand eindeutig zu charakterisieren. Eine weitere Observable $B$, die mit $A$ kommutiert, kann diese Entartung auflösen.

\begin{definitionbox}
Ein vollständiger Satz kommutierender Observablen, kurz VSKO, ist eine Menge von Observablen
\[
    A_1,A_2,\dots,A_n,
\]
die paarweise kommutieren und deren gemeinsame Eigenwerte die gemeinsamen Eigenzustände eindeutig festlegen.
\end{definitionbox}

Beispielhaft kann ein Hamiltonoperator $H$ mit einer weiteren Observable $B$ kommutieren. Wenn $H$ entartet ist, kann $B$ helfen, die Zustände innerhalb der entarteten Energieniveaus zu unterscheiden. Falls aber auch $H$ und $B$ zusammen noch Entartungen übrig lassen, sind sie kein vollständiger Satz.

\section{Drehimpuls-Kommutatoren}

\subsection{Definition des Bahndrehimpulses}

Der Bahndrehimpuls ist
\[
    \vec L=\vec X\times\vec P.
\]
In Komponenten schreibt man
\[
    L_j=\sum_{k,\ell=1}^3\varepsilon_{jk\ell}X_kP_\ell,
\]
wobei $\varepsilon_{jk\ell}$ das Levi-Civita-Symbol ist.

\begin{notebox}
Die Definition sieht fast genauso aus wie klassisch. Der Unterschied ist, dass $X_k$ und $P_\ell$ Operatoren sind. Deshalb entstehen beim Vertauschen Kommutatoren.
\end{notebox}

\subsection{Drehimpuls mit Ort und Impuls}

Die Drehimpulskomponenten erzeugen Rotationen. Mathematisch sieht man das an ihren Kommutatoren mit Ort und Impuls:
\[
    [L_j,X_k]=i\hbar\sum_{\ell=1}^3\varepsilon_{jk\ell}X_\ell,
\]
\[
    [L_j,P_k]=i\hbar\sum_{\ell=1}^3\varepsilon_{jk\ell}P_\ell.
\]

\begin{examplebox}
Für $L_z=L_3=XP_y-YP_x$ gilt
\[
\begin{aligned}
    [L_z,X]
    &=[XP_y-YP_x,X] \\
    &=X[P_y,X]+[X,X]P_y-Y[P_x,X]-[Y,X]P_x \\
    &=0+0-Y(-i\hbar)-0 \\
    &=i\hbar Y.
\end{aligned}
\]
Genauso erhält man
\[
    [L_z,Y]=-i\hbar X.
\]
Das entspricht einer infinitesimalen Rotation in der $xy$-Ebene.
\end{examplebox}

\subsection{Drehimpulsalgebra}

Die Drehimpulskomponenten kommutieren nicht miteinander. Es gilt
\[
    [L_j,L_k]=i\hbar\sum_{\ell=1}^3\varepsilon_{jk\ell}L_\ell.
\]
Konkret:
\[
    [L_x,L_y]=i\hbar L_z,
    \qquad
    [L_y,L_z]=i\hbar L_x,
    \qquad
    [L_z,L_x]=i\hbar L_y.
\]

\begin{formulabox}
Für den Gesamtdrehimpuls
\[
    L^2=L_x^2+L_y^2+L_z^2
\]
gilt
\[
    [L^2,L_j]=0
\]
für $j=1,2,3$.
\end{formulabox}

Das ist physikalisch entscheidend: Man kann den Betrag des Drehimpulses, also $L^2$, und eine Komponente, meistens $L_z$, gleichzeitig scharf messen. Man kann aber nicht alle drei Komponenten $L_x,L_y,L_z$ gleichzeitig scharf messen, weil sie nicht miteinander kommutieren.

\section{Zentralpotentiale und Drehimpulserhaltung}

Betrachten wir den Hamiltonoperator
\[
    H=\frac{\vec P^{\,2}}{2m}+V(r),
    \qquad
    r=\sqrt{X_1^2+X_2^2+X_3^2}.
\]
Das Potential hängt nur vom Abstand zum Ursprung ab. Es ist also rotationssymmetrisch.

\begin{conceptbox}
Eine Symmetrie des Hamiltonoperators führt zu einer Erhaltungsgröße. Rotationssymmetrie bedeutet Drehimpulserhaltung.
\end{conceptbox}

Für ein Zentralpotential gilt
\[
    [H,L_j]=0
\]
für alle $j$. Daher gilt auch
\[
    [H,L^2]=0.
\]
Damit können Eigenzustände von $H$ gleichzeitig als Eigenzustände von $L^2$ und $L_z$ gewählt werden. Das ist die Grundlage der Behandlung von Kugelkoordinaten, Kugelflächenfunktionen und Zentralpotentialen.

\begin{examtipbox}
Wenn in einer Aufgabe $V(r)$ gegeben ist, solltest du sofort an Rotationssymmetrie denken. Typische Konsequenzen sind:
\[
    [H,L^2]=0,
    \qquad
    [H,L_z]=0,
\]
und Zustände können mit Quantenzahlen $l,m$ klassifiziert werden.
\end{examtipbox}

\section{Varianz und Unschärfe}

\subsection{Erwartungswert und Varianz}

Sei $A$ eine Observable und $\psi$ ein normierter Zustand. Der Erwartungswert ist
\[
    \langle A\rangle=\langle\psi,A\psi\rangle.
\]
Die Varianz ist
\[
    (\Delta A)^2=\langle A^2\rangle-\langle A\rangle^2.
\]
Die Standardabweichung oder Unschärfe ist
\[
    \Delta A=\sqrt{\langle A^2\rangle-\langle A\rangle^2}.
\]

\begin{definitionbox}
Für eine Observable $A$ definiert man den zentrierten Operator
\[
    A' := A-\langle A\rangle.
\]
Dann gilt
\[
    (\Delta A)^2=\langle (A')^2\rangle=\norm{A'\psi}^2.
\]
\end{definitionbox}

\begin{derivationbox}
Da $\langle A\rangle$ eine Zahl ist, gilt
\[
\begin{aligned}
    \langle(A-\langle A\rangle)^2\rangle
    &=\langle A^2-2\langle A\rangle A+\langle A\rangle^2\rangle \\
    &=\langle A^2\rangle-2\langle A\rangle^2+\langle A\rangle^2 \\
    &=\langle A^2\rangle-\langle A\rangle^2.
\end{aligned}
\]
Für hermitesches $A$ ist $A'$ ebenfalls hermitesch, daher
\[
    \langle(A')^2\rangle=\langle A'\psi,A'\psi\rangle=\norm{A'\psi}^2\ge0.
\]
\end{derivationbox}

\subsection{Wann ist die Unschärfe null?}

Wenn $\Delta A=0$, dann gilt
\[
    \norm{(A-\langle A\rangle)\psi}^2=0.
\]
Also
\[
    (A-\langle A\rangle)\psi=0
\]
und damit
\[
    A\psi=\langle A\rangle\psi.
\]
Der Zustand ist also Eigenzustand von $A$.

\begin{conceptbox}
Eine Observable ist genau dann in einem Zustand scharf bestimmt, wenn der Zustand ein Eigenzustand der Observable ist. Dann ist die Varianz null.
\end{conceptbox}

\section{Allgemeine Unschärferelation}

\subsection{Aussage der Robertson-Unschärferelation}

\begin{formulabox}
Für zwei hermitesche Operatoren $A$ und $B$ gilt in jedem normierten Zustand
\[
    \Delta A\,\Delta B
    \ge
    \frac12\left|\langle[A,B]\rangle\right|.
\]
Diese Beziehung heißt Robertson-Unschärferelation. Für Ort und Impuls wird sie zur Heisenbergschen Unschärferelation.
\end{formulabox}

Die allgemeine Aussage ist: Wenn der Erwartungswert des Kommutators ungleich null ist, dann können beide Unschärfen nicht gleichzeitig beliebig klein sein.

\subsection{Beweis mit der Cauchy-Schwarz-Ungleichung}

Wir setzen
\[
    A'=A-\langle A\rangle,
    \qquad
    B'=B-\langle B\rangle.
\]
Dann ist
\[
    \Delta A=\norm{A'\psi},
    \qquad
    \Delta B=\norm{B'\psi}.
\]
Die Cauchy-Schwarz-Ungleichung liefert
\[
    \norm{A'\psi}^2\norm{B'\psi}^2
    \ge
    |\langle A'\psi,B'\psi\rangle|^2.
\]
Also
\[
    (\Delta A)^2(\Delta B)^2
    \ge
    |\langle A'\psi,B'\psi\rangle|^2.
\]

Jetzt ist
\[
    \langle A'\psi,B'\psi\rangle
    =\langle\psi,A'B'\psi\rangle
    =\langle A'B'\rangle.
\]
Wir zerlegen
\[
    A'B'=\frac12\{A',B'\}+\frac12[A',B'].
\]
Für hermitesche $A'$ und $B'$ ist der Erwartungswert des Antikommutators reell und der Erwartungswert des Kommutators rein imaginär. Daher ist
\[
    |\langle A'B'\rangle|^2
    \ge
    \left|\frac12\langle[A',B']\rangle\right|^2.
\]
Da Konstanten mit allem kommutieren, gilt
\[
    [A',B']=[A,B].
\]
Somit folgt
\[
    (\Delta A)^2(\Delta B)^2
    \ge
    \frac14|\langle[A,B]\rangle|^2.
\]
Wurzelziehen liefert
\[
    \Delta A\,\Delta B\ge\frac12|\langle[A,B]\rangle|.
\]

\begin{notebox}
Die noch stärkere Robertson-Schrödinger-Relation enthält zusätzlich den Antikommutatoranteil. Für viele Standardaufgaben reicht aber die schwächere, aber sehr wichtige Form
\[
    \Delta A\,\Delta B\ge\frac12|\langle[A,B]\rangle|.
\]
\end{notebox}

\section{Heisenbergsche Orts-Impuls-Unschärferelation}

Setzen wir
\[
    A=X,
    \qquad
    B=P,
\]
dann ist
\[
    [X,P]=i\hbar.
\]
Damit gilt
\[
    \langle[X,P]\rangle=i\hbar.
\]
Also folgt
\[
    \Delta X\,\Delta P
    \ge
    \frac12|i\hbar|.
\]
Da $|i\hbar|=\hbar$, erhält man
\[
    \Delta X\,\Delta P\ge\frac{\hbar}{2}.
\]

\begin{formulabox}
Die Heisenbergsche Unschärferelation für Ort und Impuls lautet
\[
    \Delta X\,\Delta P\ge\frac{\hbar}{2}.
\]
In drei Dimensionen gilt komponentenweise
\[
    \Delta X_j\,\Delta P_j\ge\frac{\hbar}{2}.
\]
Für verschiedene Komponenten gilt dagegen wegen $[X_j,P_k]=0$ für $j\ne k$ keine solche untere Schranke aus diesem Kommutator.
\end{formulabox}

\begin{conceptbox}
Die Unschärferelation ist keine Aussage über Messfehler eines schlechten Messgeräts. Sie ist eine Aussage über die Struktur quantenmechanischer Zustände: Es gibt keinen Zustand, der gleichzeitig beliebig scharf in Ort und Impuls lokalisiert ist.
\end{conceptbox}

\section{Beweis der Orts-Impuls-Unschärfe mit quadratischer Positivität}

In vielen Übungsaufgaben wird die Heisenbergsche Unschärferelation direkt über einen Hilfszustand gezeigt. Sei
\[
    X'=X-\langle X\rangle,
    \qquad
    P'=P-\langle P\rangle.
\]
Für einen reellen Parameter $\lambda$ betrachten wir
\[
    |\varphi\rangle=(X'+i\lambda P')|\psi\rangle.
\]
Da Normen nichtnegativ sind, gilt
\[
    \langle\varphi|\varphi\rangle\ge0
\]
für alle $\lambda\in\mathbb R$.

\begin{derivationbox}
Wir berechnen die Norm:
\[
\begin{aligned}
    \langle\varphi|\varphi\rangle
    &=\langle\psi|(X'-i\lambda P')(X'+i\lambda P')|\psi\rangle \\
    &=\langle (X')^2\rangle
      +i\lambda\langle X'P'\rangle
      -i\lambda\langle P'X'\rangle
      +\lambda^2\langle (P')^2\rangle \\
    &=(\Delta X)^2+\\lambda^2(\Delta P)^2+i\lambda\langle[X',P']\rangle.
\end{aligned}
\]
Da
\[
    [X',P']=[X,P]=i\hbar
\]
gilt, folgt
\[
    i\lambda\langle[X',P']\rangle
    =i\lambda i\hbar
    =-\lambda\hbar.
\]
Damit ist
\[
    \langle\varphi|\varphi\rangle
    =(\Delta X)^2+\lambda^2(\Delta P)^2-\lambda\hbar.
\]
Dieser Ausdruck muss für alle reellen $\lambda$ nichtnegativ sein:
\[
    (\Delta P)^2\lambda^2-\hbar\lambda+(\Delta X)^2\ge0.
\]
Ein quadratisches Polynom mit positivem Leitkoeffizienten ist genau dann für alle $\lambda$ nichtnegativ, wenn seine Diskriminante nicht positiv ist:
\[
    \hbar^2-4(\Delta P)^2(\Delta X)^2\le0.
\]
Daraus folgt
\[
    \Delta X\Delta P\ge\frac{\hbar}{2}.
\]
\end{derivationbox}

\begin{examtipbox}
Dieser Beweis ist klausurtypisch. Der zentrale Trick ist, einen Zustand der Form
\[
    (X'+i\lambda P')|\psi\rangle
\]
zu betrachten und seine Norm als quadratisches Polynom in $\lambda$ auszuwerten.
\end{examtipbox}

\section{Zustände minimaler Unschärfe}

\subsection{Bedingung für Gleichheit}

Die Heisenbergsche Schranke
\[
    \Delta X\Delta P\ge\frac{\hbar}{2}
\]
wird genau dann erreicht, wenn die verwendete Norm für ein bestimmtes $\lambda_0$ verschwindet:
\[
    (X'+i\lambda_0P')|\psi\rangle=0.
\]
Aus der quadratischen Minimierung ergibt sich
\[
    \lambda_0=\frac{\hbar}{2(\Delta P)^2}
    =\frac{2(\Delta X)^2}{\hbar}.
\]

\begin{formulabox}
Ein Zustand minimaler Orts-Impuls-Unschärfe erfüllt
\[
    (X-\langle X\rangle+i\lambda_0(P-\langle P\rangle))\psi=0,
\]
mit
\[
    \lambda_0=\frac{\hbar}{2(\Delta P)^2}
    =\frac{2(\Delta X)^2}{\hbar}.
\]
\end{formulabox}

\subsection{Lösung in Ortsdarstellung}

Wir lösen die Gleichung
\[
    (X'+i\lambda_0P')\psi=0
\]
in Ortsdarstellung. Setze
\[
    x_0=\langle X\rangle,
    \qquad
    p_0=\langle P\rangle.
\]
Dann ist
\[
    X'=x-x_0,
    \qquad
    P'=-i\hbar\frac{\dd}{\dd x}-p_0.
\]
Die Gleichung lautet
\[
    \left[(x-x_0)+i\lambda_0\left(-i\hbar\frac{\dd}{\dd x}-p_0\right)\right]\psi(x)=0.
\]
Also
\[
    \left[(x-x_0)+\lambda_0\hbar\frac{\dd}{\dd x}-i\lambda_0p_0\right]\psi(x)=0.
\]
Damit
\[
    \lambda_0\hbar\frac{\dd\psi}{\dd x}
    =-\left(x-x_0-i\lambda_0p_0\right)\psi.
\]
Also
\[
    \frac{1}{\psi}\frac{\dd\psi}{\dd x}
    =-\frac{x-x_0}{\lambda_0\hbar}+i\frac{p_0}{\hbar}.
\]
Integration liefert
\[
    \ln\psi(x)
    =-\frac{(x-x_0)^2}{2\lambda_0\hbar}+i\frac{p_0x}{\hbar}+C.
\]
Somit ist
\[
    \psi(x)=N\exp\left(-\frac{(x-x_0)^2}{2\lambda_0\hbar}\right)
    \exp\left(i\frac{p_0x}{\hbar}\right).
\]

\begin{formulabox}
Die Zustände minimaler Orts-Impuls-Unschärfe sind Gaußfunktionen:
\[
    \psi(x)=N\exp\left(-\frac{(x-x_0)^2}{4(\Delta X)^2}\right)
    \exp\left(i\frac{p_0x}{\hbar}\right),
\]
mit
\[
    \Delta X\Delta P=\frac{\hbar}{2}.
\]
\end{formulabox}

Hier wurde verwendet, dass
\[
    \lambda_0=\frac{2(\Delta X)^2}{\hbar}
\]
gilt. Dann ist
\[
    2\lambda_0\hbar=4(\Delta X)^2.
\]

\begin{conceptbox}
Gaußsche Wellenpakete sind keine zufällige bequeme Wahl. Sie sind genau die Zustände, die Ort und Impuls so gut wie quantenmechanisch möglich gleichzeitig lokalisieren.
\end{conceptbox}

\section{Typische Rechenaufgaben}

\subsection{Aufgabe: Kommutator von $H$ und $X_j$}

\begin{examplebox}
Gegeben sei
\[
    H=\sum_{k=1}^3\frac{P_k^2}{2m}+V(\vec X).
\]
Berechne $[H,X_j]$.
\end{examplebox}

\begin{calculationbox}
Da $V(\vec X)$ mit $X_j$ kommutiert, gilt
\[
    [H,X_j]=\sum_{k=1}^3\frac{1}{2m}[P_k^2,X_j].
\]
Mit
\[
    [P_k,X_j]=-i\hbar\delta_{jk}
\]
folgt
\[
\begin{aligned}
    [P_k^2,X_j]
    &=P_k[P_k,X_j]+[P_k,X_j]P_k \\
    &=P_k(-i\hbar\delta_{jk})+(-i\hbar\delta_{jk})P_k \\
    &=-2i\hbar\delta_{jk}P_k.
\end{aligned}
\]
Damit
\[
    [H,X_j]
    =\sum_{k=1}^3\frac{1}{2m}(-2i\hbar\delta_{jk}P_k)
    =-\frac{i\hbar}{m}P_j.
\]
\end{calculationbox}

\subsection{Aufgabe: Kommutator von $H$ und $P_j$}

\begin{examplebox}
Gegeben sei wieder
\[
    H=\sum_{k=1}^3\frac{P_k^2}{2m}+V(\vec X).
\]
Berechne $[H,P_j]$.
\end{examplebox}

\begin{calculationbox}
Da alle Impulskomponenten miteinander kommutieren, gilt
\[
    \left[\sum_{k=1}^3\frac{P_k^2}{2m},P_j\right]=0.
\]
Also
\[
    [H,P_j]=[V(\vec X),P_j].
\]
Aus
\[
    [P_j,V(\vec X)]=-i\hbar\frac{\partial V}{\partial x_j}
\]
folgt
\[
    [V(\vec X),P_j]=i\hbar\frac{\partial V}{\partial x_j}.
\]
Daher
\[
    [H,P_j]=i\hbar\frac{\partial V}{\partial x_j}.
\]
\end{calculationbox}

\subsection{Aufgabe: Heisenberg-Unschärfe prüfen}

\begin{examplebox}
Ein Zustand habe
\[
    \Delta x=\frac{1}{\sqrt{2}\rho},
    \qquad
    \Delta p=\hbar\rho.
\]
Prüfe die Heisenbergsche Unschärferelation.
\end{examplebox}

\begin{calculationbox}
Das Produkt ist
\[
    \Delta x\Delta p
    =\frac{1}{\sqrt{2}\rho}\cdot\hbar\rho
    =\frac{\hbar}{\sqrt{2}}.
\]
Da
\[
    \frac{1}{\sqrt{2}}>\frac12
\]
gilt,
\[
    \Delta x\Delta p=\frac{\hbar}{\sqrt{2}}>\frac{\hbar}{2}.
\]
Die Heisenbergsche Unschärferelation ist erfüllt. Das Minimum wird aber nicht erreicht.
\end{calculationbox}

\begin{notebox}
Wenn das Produkt größer als $\hbar/2$ ist, ist der Zustand erlaubt. Nur Produkte kleiner als $\hbar/2$ sind unmöglich. Das Minimum wird nur von speziellen Gaußzuständen erreicht.
\end{notebox}

\section{Häufige Fehler}

\begin{itemize}
    \item Man vergisst, dass $[A,B]=-[B,A]$ gilt. Das Vorzeichen ist bei $[X,P]$ und $[P,X]$ entscheidend.
    \item Man schreibt $[X,P]=\hbar$ statt $[X,P]=i\hbar$. Der Faktor $i$ ist wichtig, weil der Kommutator zweier hermitescher Operatoren antihermitesch ist.
    \item Man behandelt Operatoren wie gewöhnliche Zahlen und vertauscht sie ohne Prüfung.
    \item Man verwechselt $\Delta A$ mit $\langle A\rangle$. Die Unschärfe ist eine Standardabweichung, kein Mittelwert.
    \item Man interpretiert die Unschärferelation als technischen Messfehler statt als Zustandsaussage.
    \item Man schließt aus $[A,B]\ne0$, dass es gar keine gemeinsamen Eigenzustände geben kann. Richtig ist: Im Allgemeinen gibt es keine vollständige gemeinsame Eigenbasis. Einzelne gemeinsame Eigenzustände können trotzdem vorkommen.
\end{itemize}

\section{Zusammenfassung}

\begin{conceptbox}
Die wichtigsten Punkte dieses Kapitels sind:
\begin{itemize}
    \item Der Kommutator ist $[A,B]=AB-BA$.
    \item Ort und Impuls erfüllen $[X_j,P_k]=i\hbar\delta_{jk}$.
    \item Aus den kanonischen Kommutatoren folgen Rechenregeln wie
    \[
        [X_j,G(P_j)]=i\hbar\frac{\partial G}{\partial P_j},
        \qquad
        [P_j,F(X_j)]=-i\hbar\frac{\partial F}{\partial X_j}.
    \]
    \item Wenn $[H,A]=0$ und $A$ nicht explizit zeitabhängig ist, ist $A$ erhalten.
    \item Aus $[H,X_j]$ und $[H,P_j]$ folgen die Ehrenfest-Gleichungen.
    \item Kommutierende Observablen können simultan scharf messbar sein und gemeinsame Eigenbasen besitzen.
    \item Nichtverschwindende Kommutatoren erzwingen Unschärferelationen.
    \item Allgemein gilt $\Delta A\Delta B\ge\frac12|\langle[A,B]\rangle|$.
    \item Für Ort und Impuls gilt $\Delta X\Delta P\ge\hbar/2$.
    \item Minimale Orts-Impuls-Unschärfe wird durch Gaußzustände erreicht.
\end{itemize}
\end{conceptbox}

\section{Was du für Aufgaben sicher können solltest}

Nach diesem Kapitel solltest du folgende Standardaufgaben lösen können:
\begin{itemize}
    \item $[X_j,P_k]$ direkt in Ortsdarstellung berechnen.
    \item Produktregeln für Kommutatoren anwenden.
    \item $[X,P^n]$ und $[P,X^n]$ berechnen.
    \item $[X,G(P)]$ und $[P,F(X)]$ herleiten.
    \item $[H,X_j]$ und $[H,P_j]$ für $H=P^2/2m+V(X)$ berechnen.
    \item Das Ehrenfest-Theorem aus Kommutatoren erklären.
    \item Aus $[H,A]=0$ auf eine Erhaltungsgröße schließen.
    \item Die Heisenbergsche Unschärferelation beweisen.
    \item Zustände minimaler Unschärfe als Gaußfunktionen herleiten.
    \item Für Drehimpulse die Relationen $[L_i,L_j]=i\hbar\varepsilon_{ijk}L_k$ verwenden.
\end{itemize}
